
\subsubsection{6.1 Why SE Wins on TriviaQA}

The +0.155 AUROC gap is directly explained by the mechanism measured in RQ2. Token entropy aggregates surface variation within semantically equivalent paraphrase groups --- treating different word choices for the same wrong answer as independent uncertainty signals. Semantic entropy's NLI clustering removes this variation before computing entropy. The scale of the effect (intra-cluster TE variance = 7.152 $nats^2, 71\times$ threshold, across 76/98 questions) makes the explanation unambiguous: paraphrase noise is the dominant component of TE's within-cluster signal. This extends Huang et al. [2023]'s empirical observation that sampling-based methods outperform single-pass entropy by providing the feature-level mechanism.

The high mean cluster count (7.31 out of 10 per question) confirms that NLI clustering is actively partitioning the samples --- incorrect answers to TriviaQA questions take many different lexical forms, which SE absorbs into distinct clusters while TE counts each surface variant as an independent uncertainty signal.

\subsubsection{6.2 Why SE Loses on TruthfulQA}

TruthfulQA's adversarial misconceptions produce low-diversity model outputs: across K=10 samples, Llama-2-7B consistently generates the same wrong answer with high confidence. All K samples fall into one NLI cluster, yielding near-zero SE for incorrect answers --- making them indistinguishable from correct answers in SE's feature space. Token entropy succeeds here: a narrow, peaked token distribution (low entropy) for a deterministically wrong output is a reliable incorrectness indicator, since low TE means the model is confident in its (wrong) answer.

Three competing explanations exist for the reversal: (1) task-structure dependence --- incorrect TruthfulQA outputs are low-diversity (primary, mechanistically consistent with H-M1's paraphrase-noise account); (2) NLI model size --- H-C1 used \texttt{nli-deberta-v3-small} while H-E1 used \texttt{nli-deberta-v3-large}, a genuine confound that may inflate clustering errors on TruthfulQA's longer answer strings; (3) EM label noise --- yes/no prefix matching may be noisier than TriviaQA alias list matching (low plausibility: would deflate all AUROCs equally). The task-structure and NLI model confounds cannot be separated without the follow-up experiment described in Section 7.

\subsubsection{6.3 SCG BERTScore as a Short-QA Failure Mode}

The |SCG $-$ SE| gap of 0.336 (corrected) establishes that BERTScore-based consistency is not an SE substitute on short factual QA. BERTScore measures contextual embedding similarity, not semantic entailment --- on 1-3 word TriviaQA answers, this distinction is critical. Two answers such as ''Paris'' and ''the French capital'' may be NLI-entailed and would be grouped into the same SE cluster, but BERTScore assigns them low similarity due to surface-form divergence. Practitioners seeking lower-compute SE alternatives for short-answer tasks should use SelfCheckNLI rather than the BERTScore variant; whether NLI-based SCG recovers SE-equivalent performance is an open question \citep{Manakul2023SelfCheckGPT}.

\subsubsection{6.4 VC Degeneracy as Meta-Cognitive Failure}

The ECE = 0.430 and 5-value degenerate distribution confirm that 7B-scale instruction-tuned models lack reliable meta-cognitive access to uncertainty. A model reporting 95\% confidence on 60\% of questions while achieving 47\% accuracy is not simply underconfident or overconfident in a uniform way --- it is producing a near-constant signal that cannot be used for reliable selective abstention. This independently replicates Xiong et al. [2023] with a different elicitation prompt, confirming the finding is robust to prompt variation.

The AUROC gate failure is itself informative: VC (0.446) marginally exceeds TE (0.438) near chance, which is not evidence of VC reliability but rather evidence that both methods fail near-equally on TriviaQA at 7B scale. The VC signal retains just enough variance from the two questions where the model reported 0\% confidence (both incorrect) to match TE by AUROC --- this is not a generalizable uncertainty estimation strategy.

\subsubsection{6.5 Limitations}

\textbf{L1 (Task-structure scope):} SE > TE holds on TriviaQA but not TruthfulQA. The task-structure condition --- paraphrase-rich vs. deterministic-wrong incorrect outputs --- is identified as a contribution, not merely a failure. Identifying \textit{when} a method works is more informative than an unconditional ranking claim.

\textbf{L2 (N=98 pilot statistical power):} CI half-widths of approximately 0.05--0.07 provide ample power for the +0.155 primary finding. TruthfulQA results (AUROC range 0.44--0.51, N=141) have overlapping CIs and should be interpreted as directional. The pre-specified N=500 extension protocol would be required for definitive TruthfulQA ranking.

\textbf{L3 (NLI model size confound in H-C1):} H-C1 changes both the benchmark and the NLI model size simultaneously (\texttt{nli-deberta-v3-small} vs. \texttt{nli-deberta-v3-large}). The TruthfulQA reversal cannot be attributed purely to task structure without a follow-up experiment holding the NLI model fixed.

\textbf{L4 (Sign convention inconsistency across pipelines):} H-M3 and H-M4 report SE AUROC = 0.286 (uninverted orientation). All within-pipeline comparisons are internally consistent; corrected values (SE = 0.714) are provided. The P2 refutation (SCG $\neq$ SE) is robust to sign conventions: the raw delta (0.092) and corrected delta (0.336) both exceed the 0.03 gate by $3\times$ and $11\times$ respectively.

\textbf{L5 (H-M2 ablation design flaw):} The planned NLI-clustering ablation was structurally invalidated because the ablated predictor is a monotone transform of SE, making AUROC delta = 0 by construction. Mechanistic evidence relies on H-M1's direct intra-cluster variance measurement.

\subsubsection{6.6 Broader Implications}

This work characterizes when SE is and is not a reliable uncertainty estimator, providing practitioners with empirical guidance for method selection. The task-structure condition --- paraphrase-rich vs. deterministic-wrong incorrect outputs --- is a benchmark property that can be evaluated before committing to a method. The VC degeneracy finding has direct safety implications: 7B-scale models expressing 95\% confidence while answering incorrectly approximately 53\% of the time should not be used as uncertainty proxies without calibration correction.

